\documentclass[a4paper,12pt]{article}

\usepackage[T2A]{fontenc}
\usepackage[utf8]{inputenc}
\usepackage[russian]{babel}

\usepackage{amsmath,amssymb,mathtools}
\usepackage{helvet}
\renewcommand{\familydefault}{\sfdefault}

\usepackage{icomma}
\usepackage[a4paper,margin=20mm]{geometry}
\usepackage{microtype}
\usepackage{tikz}
\usetikzlibrary{calc,arrows.meta,positioning,shapes.geometric}

\usepackage{tabularx,booktabs}
\usepackage{enumitem}
\usepackage{hyperref}
\usepackage{parskip}
\usepackage{fancyhdr}
\usepackage{setspace}
\usepackage{xcolor}

\definecolor{accent}{HTML}{008080}
\definecolor{darkaccent}{HTML}{004D4D}
\definecolor{textgray}{HTML}{333333}
\definecolor{softgray}{HTML}{E9EEEE}

\hypersetup{hidelinks}

\geometry{
  left=20mm,
  right=20mm,
  top=20mm,
  bottom=22mm
}

\onehalfspacing
\setlength{\parindent}{0pt}
\setlength{\parskip}{5pt}

\setlist[itemize]{
  leftmargin=7mm,
  itemsep=2pt,
  topsep=3pt
}

\setlist[enumerate]{
  leftmargin=8mm,
  itemsep=3pt,
  topsep=3pt
}

\pagestyle{fancy}
\fancyhf{}
\fancyhead[L]{\small\color{textgray}Показательно-логарифмические неравенства}
\fancyhead[R]{\small\color{textgray}ЕГЭ}
\fancyfoot[C]{\small\thepage}
\renewcommand{\headrulewidth}{0.3pt}
\renewcommand{\headrule}{%
  \hbox to\headwidth{\color{accent!45}\leaders\hrule height \headrulewidth\hfill}
}

\newcommand{\topicline}{%
  \par\noindent
  {\color{accent!55}\rule{\linewidth}{0.6pt}}
  \par
}

\tikzset{
  flowstart/.style={
    ellipse,
    draw=darkaccent,
    line width=0.9pt,
    minimum width=4.3cm,
    minimum height=1cm,
    align=center,
    font=\small
  },
  flowaction/.style={
    rectangle,
    rounded corners=3pt,
    draw=accent,
    line width=0.8pt,
    text width=8.2cm,
    minimum height=1.15cm,
    align=center,
    font=\small,
    inner sep=5pt
  },
  flowdecision/.style={
    diamond,
    aspect=2.2,
    draw=darkaccent,
    line width=0.8pt,
    text width=4.8cm,
    align=center,
    font=\small,
    inner sep=2pt
  },
  flowarrow/.style={
    -{Latex[length=3mm,width=2mm]},
    line width=0.8pt,
    draw=textgray
  },
  flowlabel/.style={
    font=\small,
    fill=white,
    inner sep=1.5pt
  }
}

\begin{document}

% =========================================================
% ТИТУЛЬНЫЙ ЛИСТ
% =========================================================

\begin{titlepage}
\thispagestyle{empty}

\begin{center}

\vspace*{24mm}

{\color{darkaccent}\Large\bfseries ЧЕК-ЛИСТ}

\vspace{8mm}

{\Huge\bfseries
Показательно-логарифмическое\\[2mm]
неравенство
}

\vspace{7mm}

{\Large
Замены, разложение на множители,\\
ОДЗ и метод интервалов
}

\vspace{16mm}

{\color{accent}\rule{0.7\linewidth}{0.8pt}}

\vspace{15mm}

{\large
ЕГЭ по профильной математике
}

\vspace{5mm}

{\large
Путь А
}

\vfill

{\large
\textbf{Лёвин Артём Александрович}\\[2mm]
эксклюзивно для Тимофея Васильченко
}

\vspace{7mm}

{\large \today}

\end{center}

\begin{tikzpicture}[remember picture,overlay]
\draw[
  accent!45,
  line width=1.3pt
]
([xshift=14mm,yshift=13mm]current page.south west)
.. controls
([xshift=62mm,yshift=27mm]current page.south west)
and
([xshift=-74mm,yshift=8mm]current page.south east)
..
([xshift=-14mm,yshift=16mm]current page.south east);
\end{tikzpicture}

\end{titlepage}

% =========================================================
% ОСНОВНОЙ ЧЕК-ЛИСТ
% =========================================================

\section*{\color{darkaccent}1. Главная идея задачи}
\topicline

Рассмотрим типичное сложное неравенство:
\[
\frac{
9^x-3^{x+2}+8
}{
\log_{\frac16}^{\,2}(5^x-2)
+
\log_{\frac16}\!\left((5^x-2)^2\right)
+1
}
\le 0.
\]

Внешний вид громоздкий, однако структура задачи состоит из нескольких стандартных действий:

\begin{enumerate}
\item привести показательные выражения к одной повторяющейся конструкции;
\item выполнить замену;
\item разложить квадратный трёхчлен на множители;
\item корректно преобразовать логарифм чётной степени;
\item выписать ОДЗ;
\item свернуть знаменатель в квадрат;
\item учесть, что квадрат знаменателя положителен только при условии, что он не равен нулю;
\item решить оставшееся произведение методом интервалов;
\item пересечь результат со всеми ограничениями.
\end{enumerate}

\section*{\color{darkaccent}2. Показательная часть: ищите повторяющееся выражение}
\topicline

Если одновременно встречаются
\[
a^{2x},\qquad a^x,
\]
естественная замена:
\[
t=a^x,\qquad t>0.
\]

В нашей задаче:
\[
9^x=(3^2)^x=3^{2x}=(3^x)^2,
\]
а также
\[
3^{x+2}=3^x\cdot3^2=9\cdot3^x.
\]

Поэтому
\[
9^x-3^{x+2}+8
=
(3^x)^2-9\cdot3^x+8.
\]

Вводим
\[
t=3^x.
\]

Получаем квадратный трёхчлен:
\[
t^2-9t+8.
\]

Его корни:
\[
t_1=1,\qquad t_2=8,
\]
поэтому
\[
t^2-9t+8=(t-1)(t-8).
\]

Возвращаем замену:
\[
\boxed{
9^x-3^{x+2}+8
=
(3^x-1)(3^x-8)
}.
\]

\subsection*{\color{accent}Триггер}

Увидели выражения вида
\[
a^{2x},\quad a^x
\]
--- попробуйте замену
\[
t=a^x.
\]

Увидели квадратный трёхчлен
\[
t^2+bt+c
\]
--- проверьте возможность разложения на множители.

\section*{\color{darkaccent}3. Степень степени}
\topicline

Ключевое правило:
\[
(a^m)^n=a^{mn}.
\]

Поэтому
\[
(3^x)^2=3^{2x}.
\]

Важно:
\[
(3^x)^2\ne 3^{x^2}.
\]

При возведении степени в степень показатели \textbf{перемножаются}.

\section*{\color{darkaccent}4. Логарифм чётной степени и модуль}
\topicline

Для
\[
a>0,\qquad a\ne1,\qquad u\ne0
\]
справедливо:
\[
\boxed{
\log_a(u^{2n})=2n\log_a|u|
}.
\]

В частности,
\[
\log_a(u^2)=2\log_a|u|.
\]

Модуль принципиален: из условия
\[
u^2>0
\]
следует только
\[
u\ne0,
\]
знак самого \(u\) заранее неизвестен.

Для нечётной степени при существовании логарифма требуется \(u>0\):
\[
\log_a(u^{2n+1})
=
(2n+1)\log_a u.
\]

В нашей задаче:
\[
\log_{\frac16}\!\left((5^x-2)^2\right)
=
2\log_{\frac16}|5^x-2|.
\]

\section*{\color{darkaccent}5. ОДЗ снимает модуль}
\topicline

В знаменателе присутствует
\[
\log_{\frac16}(5^x-2).
\]

Следовательно,
\[
5^x-2>0.
\]

Отсюда
\[
5^x>2,
\]
а поскольку \(5>1\),
\[
\boxed{x>\log_5 2}.
\]

На ОДЗ известно:
\[
5^x-2>0.
\]

Поэтому
\[
|5^x-2|=5^x-2.
\]

Значит,
\[
\log_{\frac16}\!\left((5^x-2)^2\right)
=
2\log_{\frac16}(5^x-2).
\]

\section*{\color{darkaccent}6. Сворачиваем знаменатель}
\topicline

Обозначим
\[
t=\log_{\frac16}(5^x-2).
\]

Тогда знаменатель превращается в
\[
t^2+2t+1.
\]

Используем формулу сокращённого умножения:
\[
t^2+2t+1=(t+1)^2.
\]

Следовательно,
\[
\boxed{
\log_{\frac16}^{\,2}(5^x-2)
+
\log_{\frac16}\!\left((5^x-2)^2\right)
+1
=
\left(
\log_{\frac16}(5^x-2)+1
\right)^2
}.
\]

\subsection*{\color{accent}Триггер}

Увидели
\[
A^2+2A+1
\]
--- сразу проверяйте формулу
\[
(A+1)^2.
\]

Формулы сокращённого умножения полезно узнавать в обе стороны:
\[
(A+1)^2
\longleftrightarrow
A^2+2A+1.
\]

\section*{\color{darkaccent}7. Квадрат знаменателя: два разных утверждения}
\topicline

После преобразований исходная дробь имеет вид
\[
\frac{
(3^x-1)(3^x-8)
}{
\left(
\log_{\frac16}(5^x-2)+1
\right)^2
}
\le0.
\]

Для любого допустимого \(x\)
\[
\left(
\log_{\frac16}(5^x-2)+1
\right)^2\ge0.
\]

Однако знаменатель дроби не имеет права обращаться в нуль.

Поэтому обязательно записываем:
\[
\log_{\frac16}(5^x-2)+1\ne0.
\]

Решим запрещённое равенство:
\[
\log_{\frac16}(5^x-2)=-1.
\]

По определению логарифма:
\[
5^x-2=
\left(\frac16\right)^{-1}=6.
\]

Следовательно,
\[
5^x=8,
\]
то есть
\[
\boxed{x\ne\log_5 8}.
\]

Итак, на области определения знаменатель строго положителен:
\[
\left(
\log_{\frac16}(5^x-2)+1
\right)^2>0.
\]

Следовательно, знак всей дроби определяется числителем.

\section*{\color{darkaccent}8. Получаем систему}
\topicline

Исходное неравенство равносильно системе
\[
\begin{cases}
(3^x-1)(3^x-8)\le0,\\[1mm]
x>\log_5 2,\\[1mm]
x\ne\log_5 8.
\end{cases}
\]

Такой переход допустим именно потому, что знаменатель на допустимых значениях \(x\) положителен.

\section*{\color{darkaccent}9. Решение показательных уравнений}
\topicline

Для метода интервалов найдём нули множителей.

Первый:
\[
3^x-1=0.
\]

Отсюда
\[
3^x=1=3^0,
\]
следовательно,
\[
x=0.
\]

Второй:
\[
3^x-8=0,
\]
то есть
\[
3^x=8.
\]

\subsection*{\color{accent}Способ 1. Логарифмирование}

Применяем логарифм по основанию \(3\):
\[
\log_3(3^x)=\log_3 8.
\]

Используя
\[
\log_a(a^x)=x,
\]
получаем
\[
\boxed{x=\log_3 8}.
\]

\subsection*{\color{accent}Способ 2. Основное логарифмическое тождество}

При
\[
a>0,\qquad a\ne1,\qquad b>0
\]
справедливо:
\[
\boxed{a^{\log_a b}=b}.
\]

Поэтому
\[
8=3^{\log_3 8}.
\]

Тогда
\[
3^x=3^{\log_3 8},
\]
откуда
\[
x=\log_3 8.
\]

Первый способ универсален и особенно полезен в более сложных показательных уравнениях.

\section*{\color{darkaccent}10. Метод интервалов}
\topicline

Решаем
\[
(3^x-1)(3^x-8)\le0.
\]

Критические точки:
\[
x=0,
\qquad
x=\log_3 8.
\]

Так как функция \(3^x\) возрастает,
\[
0<\log_3 8.
\]

Знаки произведения:

\begin{center}
\begin{tikzpicture}[x=1cm,y=1cm]
\draw[-{Latex[length=2.5mm]},line width=0.8pt]
  (0,0)--(11,0);

\draw[line width=0.8pt] (3,-0.12)--(3,0.12);
\draw[line width=0.8pt] (8,-0.12)--(8,0.12);

\filldraw[fill=white,draw=darkaccent,line width=0.9pt]
  (3,0) circle (2.1pt);
\filldraw[fill=white,draw=darkaccent,line width=0.9pt]
  (8,0) circle (2.1pt);

\node[below=4pt] at (3,0) {$0$};
\node[below=4pt] at (8,0) {$\log_3 8$};

\node[above=5pt] at (1.5,0) {$+$};
\node[above=5pt] at (5.5,0) {$-$};
\node[above=5pt] at (9.5,0) {$+$};
\end{tikzpicture}
\end{center}

Для знака «\(\le0\)» выбираем отрицательный промежуток вместе с нулями:
\[
x\in[0;\log_3 8].
\]

\section*{\color{darkaccent}11. Финальное пересечение}
\topicline

Имеем одновременно:
\[
x\in[0;\log_3 8],
\]
\[
x>\log_5 2,
\]
\[
x\ne\log_5 8.
\]

Для упорядочивания точек заметим:
\[
0<\log_5 2<\log_5 8.
\]

Кроме того,
\[
\log_5 8
=
3\log_5 2,
\qquad
\log_3 8
=
3\log_3 2.
\]

Так как
\[
5>3>1,
\]
для числа \(2>1\)
\[
\log_5 2<\log_3 2.
\]

Следовательно,
\[
\log_5 8<\log_3 8.
\]

Итак,
\[
\boxed{
x\in
(\log_5 2;\log_5 8)
\cup
(\log_5 8;\log_3 8]
}.
\]

\section*{\color{darkaccent}12. Что следует узнавать автоматически}
\topicline

\begin{tabularx}{\linewidth}{@{}p{0.35\linewidth}X@{}}
\toprule
\textbf{Что Вы видите} & \textbf{Что стоит проверить} \\
\midrule
\(a^{2x}\) и \(a^x\) &
Замена \(t=a^x\). \\[1mm]

Квадратный трёхчлен &
Разложение на множители. \\[1mm]

\(\log_a(u^{2n})\) &
Вынесение степени:
\(2n\log_a|u|\). \\[1mm]

\(A^2+2A+1\) &
Свёртка в \((A+1)^2\). \\[1mm]

Квадрат в знаменателе &
Знаменатель неотрицателен, но отдельно требуется исключить его ноль. \\[1mm]

Произведение множителей &
Метод интервалов. \\[1mm]

Уравнение \(a^x=b\) &
\(x=\log_a b\). \\[1mm]

Основание \(\frac1a\) &
Полезно помнить:
\(\log_{1/a}u=-\log_a u\). \\
\bottomrule
\end{tabularx}

\section*{\color{darkaccent}13. Типичные ошибки}
\topicline

\begin{enumerate}
\item
Писать
\[
(3^x)^2=3^{x^2}.
\]
Верно:
\[
(3^x)^2=3^{2x}.
\]

\item
Преобразовывать
\[
\log_a(u^2)
\]
как
\[
2\log_a u
\]
без проверки знака \(u\).

Корректная общая формула:
\[
\log_a(u^2)=2\log_a|u|.
\]

\item
Удалять модуль до записи ОДЗ.

В данной задаче
\[
|5^x-2|=5^x-2
\]
только после установления
\[
5^x-2>0.
\]

\item
Говорить, что квадрат знаменателя всегда положителен.

Квадрат может быть равен нулю. Для дроби это запрещено.

\item
Забывать условие
\[
x\ne\log_5 8.
\]

\item
Решить только числитель и сразу записать ответ.

Финальный этап --- пересечение с ОДЗ и всеми исключёнными значениями.

\item
При работе с логарифмическими неравенствами забывать о основании.

Если
\[
0<a<1,
\]
функция
\[
y=\log_a x
\]
убывает, поэтому при переходе от логарифмического неравенства к сравнению аргументов направление знака меняется.

\item
Сокращать множитель дроби, не сохранив условие его ненулевого значения.

Если множитель стоял в знаменателе, его ноль остаётся запрещённым даже после алгебраического сокращения.
\end{enumerate}



% =========================================================
% КРАТКИЙ ЧЕК-ЛИСТ ПЕРЕД ЭКЗАМЕНОМ
% =========================================================

\section*{\color{darkaccent}14. Проверка решения перед записью ответа}
\topicline

Перед завершением подобной задачи проверьте:

\begin{itemize}
\item[$\square$]
Выписано ОДЗ каждого логарифма.

\item[$\square$]
Учтён запрет деления на ноль.

\item[$\square$]
При снятии чётной степени под логарифмом появился модуль.

\item[$\square$]
Модуль раскрыт только после установления знака подмодульного выражения.

\item[$\square$]
При замене \(t=a^x\) учтено \(t>0\), если это существенно.

\item[$\square$]
Квадратный трёхчлен разложен без потери множителей.

\item[$\square$]
Нули числителя и знаменателя различаются по смыслу:
нули числителя могут войти в ответ,
нули знаменателя запрещены.

\item[$\square$]
Знаки на интервалах определены корректно.

\item[$\square$]
Если основание логарифма меньше \(1\),
при сравнении аргументов учтено убывание логарифмической функции.

\item[$\square$]
В конце выполнено пересечение с ОДЗ.

\item[$\square$]
Все запрещённые точки удалены из окончательного множества.
\end{itemize}

% =========================================================
% ТРЕНИРОВОЧНЫЙ БЛОК
% =========================================================

\clearpage
\section*{\color{darkaccent}Тренировочный блок --- Путь А}
\topicline

\textbf{Задача.}
Решите неравенства. Решение оформляйте последовательно:
ОДЗ \(\rightarrow\) преобразование числителя \(\rightarrow\)
преобразование знаменателя \(\rightarrow\)
метод интервалов \(\rightarrow\) пересечение ограничений.

\vspace{3mm}

\begin{enumerate}

\item
\[
\frac{
4^x-5\cdot2^x+4
}{
\log_{\frac13}^{\,2}(3^x-1)
+
\log_{\frac13}\!\left((3^x-1)^2\right)
+1
}
\le0.
\]



\end{enumerate}



\end{document}